% !TeX root = ../../Book5.tex
% 纲要标题：### 18.4 最高权分类
% 纲要位置：工作记录/计划纲要.md，第 4304 行。
% * 支配整权分类有限维简单模
% * 存在性、唯一性与完备性
% * 经典表示及基本权的例子
% * 存在性通过 Verma 模商去简单根方向的适当幂关系构造；结合 \(\mathfrak{sl}_2\) 局部幂零性和权集控制证明有限维性，不借用下一章特征标公式倒推
% ---

\section{最高权分类}
\label{b5:sec:1804}

最高权必须支配且为整权，已由每个简单根方向的三维子代数看出。充分性需要把这些方向同时拼合起来。本节先构造一个非零商模，再用局部幂零性和 Weyl 群对称性限制其权集；整个有限维性证明不使用下一章的特征标公式。

\subsection{支配整权与有限维简单模}
\label{b5:subsec:1804:01}

设 \(\lambda\in P^+\)，记 \(m_i=\lambda(h_i)\)。在 \(M(\lambda)\) 中取向量
\[
u_i=f_i^{m_i+1}v_\lambda,
\qquad
V(\lambda)=M(\lambda)\Big/\sum_{i=1}^\ell U(\mathfrak g)u_i.
\]
这一构造只规定简单负根方向的幂关系。虽然只有有限个关系，仍须证明所得商模非零且有限维。

\begin{lemma}\label{b5:lem:integrable-verma-quotient}
设 \(\mathfrak g\) 为有限维复半单李代数，固定 Cartan 子代数 \(\mathfrak h\)、正根及标准简单根三元组 \(e_i,f_i,h_i\)（\(1\le i\le\ell\)）。设 \(\lambda\in\mathfrak h^*\) 为支配整权，记 \(m_i=\lambda(h_i)\)，\(M(\lambda)\) 为相应 Verma 模，\(v_\lambda\) 为其标准最高权生成元。则商模
\[
V(\lambda)=M(\lambda)\Big/\sum_{i=1}^\ell U(\mathfrak g)f_i^{m_i+1}v_\lambda
\]
非零，且为最高权 \(\lambda\) 的模。其中每个 \(e_i,f_i\) 都局部幂零，即对每个向量都存在一个充分高的幂将它零化。
\end{lemma}
\begin{proof}
首先，每个 \(u_i\) 都被全部 \(e_j\) 零化。对 \(j=i\)，\(\mathfrak{sl}_2\) 公式给出
\[
e_i f_i^{m_i+1}v_\lambda
=(m_i+1)\Bigl(m_i-(m_i+1)+1\Bigr)f_i^{m_i}v_\lambda=0.
\]
对 \(j\ne i\)，简单根之差不是根，所以 \([e_j,f_i]=0\)，从而 \(e_ju_i=0\)。因此 \(u_i\) 是权 \(\lambda-(m_i+1)\alpha_i\) 的最高权向量。由 PBW，它生成的子模只有低于或等于该权的权，不含 \(\lambda\)。这些子模之和仍不含最高权线，所以 \(V(\lambda)\ne0\)。

在商模的生成元 \(\bar v_\lambda\) 上，已有 \(e_i\bar v_\lambda=0\) 和 \(f_i^{m_i+1}\bar v_\lambda=0\)。根弦定理说明 \(\operatorname{ad}e_i,\operatorname{ad}f_i\) 在 \(\mathfrak g\) 上幂零。因此\cref{b5:prop:enveloping-adjoint-nilpotence}的两个假设都满足：伴随作用在李代数上局部幂零，而且算子的一次或适当次幂零化循环生成元。该命题于是给出整个 \(V(\lambda)\) 上的局部幂零性。

具体地，对 \(x=e_i\) 或 \(f_i\) 及固定 \(u\in U(\mathfrak g)\)，若
\((\operatorname{ad}x)^a(u)=0\)、\(x^b\bar v_\lambda=0\)，所引命题给出的界
\(N\ge a+b-1\) 保证 \(x^Nu\bar v_\lambda=0\)。因此这里并未假设一个统一的幂零指数零化整个无限维候选模。
\end{proof}

\begin{lemma}\label{b5:lem:integrable-weight-symmetry}
设 \(\mathfrak g\) 为有限维复半单李代数，固定 Cartan 子代数 \(\mathfrak h\)、正根及标准简单根三元组 \(e_i,f_i,h_i\)（\(1\le i\le\ell\)），以 \(W\) 表示相应的 Weyl 群。设复 \(\mathfrak g\) 模 \(V\) 具有相对于 \(\mathfrak h\) 的权分解和有限维权空间，且每个 \(e_i,f_i\) 都局部幂零。则对每个 \(i\)，每个向量都包含在一个有限维的根子代数 \(\operatorname{span}_{\mathbb C}\{e_i,f_i,h_i\}\cong\mathfrak{sl}_2(\mathbb C)\) 的子模中；对每个 \(\mu\in\mathfrak h^*\) 与 \(w\in W\)，有
\[
\dim V_\mu=\dim V_{w\mu}.
\]
\end{lemma}
\begin{proof}
先取权向量 \(v\)。第 \(i\) 个 \(\mathfrak{sl}_2\) 的 PBW 单项式可按 \(f_i^a h_i^b e_i^c\) 排序。只有有限多个 \(e_i^cv\) 非零；每个仍是权向量，所以 \(h_i^b\) 只给标量。再对每个这样的向量使用 \(f_i\) 的局部幂零性，便知 \(U\Bigl(\mathfrak{sl}_2^{(i)}\Bigr)v\) 有限维。任意向量是有限个权向量之和，所以结论对它们也成立。

设 \(r=\mu(h_i)\ge0\)。在有限维 \(\mathfrak{sl}_2\) 模中，分类说明 \(f_i^r\) 在 \(h_i\) 权为 \(r\) 的空间上单射。对 \(0\ne v\in V_\mu\)，取刚构造的有限维子模应用此结论，得 \(f_i^rv\ne0\)；其完整 \(\mathfrak h\) 权为 \(\mu-r\alpha_i=s_i\mu\)。故 \(f_i^r:V_\mu\to V_{s_i\mu}\) 单射。若 \(r\le0\)，同理用 \(e_i^{-r}\)。再从 \(s_i\mu\) 出发得到反向维数不等式，所以两者维数相等。这里 \(r\) 的整性也来自包含 \(v\) 的有限维 \(\mathfrak{sl}_2\) 模。简单反射生成 \(W\)，所以结论对所有 \(w\) 成立。
\end{proof}

\begin{lemma}\label{b5:lem:highest-integrable-finite}
设 \(\mathfrak g\) 为有限维复半单李代数，固定 Cartan 子代数 \(\mathfrak h\) 及正根，记简单根为 \(\alpha_1,\ldots,\alpha_\ell\)、对应余根为 \(h_i\)，Weyl 群为 \(W\)，令 \(Q_+=\sum_i\mathbb Z_{\ge0}\alpha_i\)。设 \(\lambda\in\mathfrak h^*\) 为支配整权。若复 \(\mathfrak g\) 权模 \(V\) 具有有限维权空间，权集包含于 \(\lambda-Q_+\)，且其权集在 \(W\) 下稳定，则 \(V\) 有限维。特别地，相应 Verma 模 \(M(\lambda)\) 的商 \(V(\lambda)=M(\lambda)/\sum_iU(\mathfrak g)f_i^{\lambda(h_i)+1}v_\lambda\) 有限维；这里 \(f_i\) 为标准简单负根向量，\(v_\lambda\) 为标准最高权生成元。
\end{lemma}
\begin{proof}
令 \(\rho^\vee=\frac12\sum_{\alpha>0}h_\alpha\)。正余根在简单余根下具有非负系数，且每个简单余根都出现，所以
\[
\rho^\vee=\sum_i c_i h_i,\qquad c_i>0.
\]
此外，简单反射置换除对应简单余根外的正余根，故
\(s_i\rho^\vee=\rho^\vee-h_i\)。与反射公式
\(s_i h=h-\alpha_i(h)h_i\) 比较，得到 \(\alpha_i(\rho^\vee)=1\)。

由\cref{b5:prop:weyl-dominant-order}，每个整权的 Weyl 群轨道有唯一支配代表。因为权集包含于 \(\lambda-Q_+\subseteq P\)，全部权均为整权。设 \(\mu\) 是 \(V\) 的一个支配权。写
\(\mu=\sum_i a_i\omega_i\)，其中在当前格中 \(a_i=\mu(h_i)\in\mathbb Z_{\ge0}\)。因为 \(\lambda-\mu\in Q_+\)，有
\[
0\le\mu(\rho^\vee)=\sum_i c_i a_i\le\lambda(\rho^\vee).
\]
每个 \(c_i>0\)，所以每个非负整数 \(a_i\) 都有上界，可能的支配权只有有限多个。每个权与其中一个支配权位于同一轨道，而 \(W\) 有限，故整个权集有限。再用每个权空间有限维，得到 \(V\) 有限维。

对 \(V(\lambda)\)，PBW 给出权集包含关系及有限维权空间，\cref{b5:lem:integrable-verma-quotient}和\cref{b5:lem:integrable-weight-symmetry}给出 Weyl 稳定性。因此所有假设均满足。
\end{proof}

\subsection{存在性、唯一性与完备性}
\label{b5:subsec:1804:02}

\begin{theorem}[最高权分类]\label{b5:thm:highest-weight-classification}
设 \(\mathfrak g\) 为有限维复半单李代数，固定 Cartan 子代数 \(\mathfrak h\) 及相应根系的一组正根。令 \(P^+\subseteq\mathfrak h^*\) 为相对于此选择的支配整权集，\(L(\lambda)\) 为相应 Verma 模 \(M(\lambda)\) 的唯一简单商。映射
\[
\lambda\longmapsto L(\lambda),\qquad\lambda\in P^+,
\]
给出支配整权与有限维简单 \(\mathfrak g\) 模同构类之间的双射。等价地，Verma 模的唯一简单商 \(L(\lambda)\) 有限维，当且仅当 \(\lambda\) 支配且为整权。
\end{theorem}
\begin{proof}
若 \(\lambda\in P^+\)，前一小节构造了非零有限维最高权模 \(V(\lambda)\)。由\cref{b5:thm:verma-simple-quotient}，它满射到 \(L(\lambda)\)，所以后者非零、有限维且简单，证明存在性。

反过来，每个有限维简单模都由\cref{b5:prop:finite-simple-highest}成为最高权模，其最高权由\cref{b5:prop:finite-highest-dominant}属于 \(P^+\)；Verma 泛性质及唯一简单商定理又说明它同构于该 \(L(\lambda)\)，证明完备性。不同 \(\lambda\) 对应的简单最高权模不同构，已在唯一简单商定理中证明，故得到单射。三步分别解决存在、穷尽和唯一的问题。
\end{proof}

结合\cref{b5:thm:weyl-complete-reducibility}，任意有限维模都可写为
\(\bigoplus_{\lambda\in P^+}L(\lambda)^{\oplus n_\lambda}\)，其中只有有限多个 \(n_\lambda\ne0\)。这些重数由 Hom 空间决定，分解中的具体子空间则通常依赖选择。

\subsection{经典表示与基本权}
\label{b5:subsec:1804:03}

对 \(\mathfrak{sl}_r\)，取 \(h_i=E_{ii}-E_{i+1,i+1}\) 和正根向量 \(E_{ij}\)、\(i<j\)。基本权由 \(\omega_i(h_j)=\delta_{ij}\) 确定。自然模 \(V=\mathbb C^r\) 的最高权为 \(\omega_1\)，其对偶最高权为 \(\omega_{r-1}\)。

\begin{proposition}\label{b5:prop:type-a-fundamental-exterior}
设 \(r\ge2\) 为整数，\(V=\mathbb C^r\) 为 \(\mathfrak{sl}_r(\mathbb C)\) 的自然模。取对角 Cartan 子代数及上三角正根，令 \(h_i=E_{ii}-E_{i+1,i+1}\)，基本权 \(\omega_i\) 由 \(\omega_i(h_j)=\delta_{ij}\)（\(1\le i,j<r\)）确定。对整数 \(1\le j<r\)，外幂 \(\bigwedge^jV\) 简单，最高权为 \(\omega_j\)。对整数 \(m\ge0\)，\(\operatorname{Sym}^mV\) 简单，最高权为 \(m\omega_1\)。
\end{proposition}
\begin{proof}
外幂的标准楔积 \(e_{i_1}\wedge\cdots\wedge e_{i_j}\)、\(i_1<\cdots<i_j\)，有互不相同的 \(\mathfrak h\) 权，每个权空间一维。向量 \(e_1\wedge\cdots\wedge e_j\) 被所有正根向量零化；对 \(h_i\) 的取值为 \(\delta_{ij}\)，所以权为 \(\omega_j\)。

若一个非零子模存在，它按权分解，故含某个标准楔积。若其指标集合不是 \(\{1,\ldots,j\}\)，选一个有空缺的较小指标和集合中较大的指标，用相应 \(E_{ab}\)、\(a<b\) 将较大指标换成空缺，得到非零楔积。反复操作使指标和下降，最终得到 \(e_1\wedge\cdots\wedge e_j\)。反向使用负根矩阵单位能从该向量产生全部标准楔积。因此任何非零子模都是整个外幂。

对称幂可用次数 \(m\) 的单项式 \(e_1^{a_1}\cdots e_r^{a_r}\) 为基；固定总次数后，它们的 \(\mathfrak{sl}_r\) 权仍互不相同。\(E_{i,i+1}\) 把一个 \(e_{i+1}\) 换成 \(e_i\)，系数为非零整数 \(a_{i+1}\)。从任一单项式反复提升可得到 \(e_1^m\)，反向下降又生成全部单项式。因此同样的权空间论证证明简单性。\(e_1^m\) 的最高权为 \(m\omega_1\)。
\end{proof}

对简单李代数，伴随模简单。取高度最大的正根及其非零根向量，正根升算子都将它零化；这个向量因简单性而生成整个伴随模。\cref{b5:prop:highest-weight-support}于是说明全部根都不高于该根，而最高权生成元的权唯一，所以得到唯一的最高根。这里使用的是伴随模的简单性，不能仅凭每个 Weyl 轨道有唯一支配代表来排除不同轨道。对 \(\mathfrak{sl}_3\)，最高根为 \(\alpha_1+\alpha_2=\omega_1+\omega_2\)，故伴随模为 \(L(\omega_1+\omega_2)\)，维数八。其非零权为六个根，零权空间为二维 Cartan 子代数。这给出一个权重数大于一的例子。

\subsection{Verma 模商、局部幂零性与有限维性}
\label{b5:subsec:1804:04}

\begin{proposition}\label{b5:prop:finite-highest-presentation}
设 \(\mathfrak g\) 为有限维复半单李代数，固定 Cartan 子代数 \(\mathfrak h\)、正根及标准简单根三元组 \(e_i,f_i,h_i\)（\(1\le i\le\ell\)）。设 \(\lambda\in\mathfrak h^*\) 为支配整权，令 \(m_i=\lambda(h_i)\)。记相应 Verma 模为 \(M(\lambda)\)，标准最高权生成元为 \(v_\lambda\)，唯一简单商为 \(L(\lambda)\)。则
\[
L(\lambda)\cong
M(\lambda)\Big/\sum_i U(\mathfrak g)f_i^{m_i+1}v_\lambda.
\]
因此有限维简单最高权模由一个向量 \(v\) 及关系
\[
hv=\lambda(h)v,\qquad e_iv=0,\qquad f_i^{m_i+1}v=0
\]
给出；这里 \(h\) 遍历 \(\mathfrak h\)，\(i\) 遍历简单根。
\end{proposition}
\begin{proof}
已经证明 \(V(\lambda)\) 非零且有限维，故由 Weyl 完全可约性分解为简单模的直和。将最高权生成元 \(\bar v_\lambda\) 投到每个简单直和项；每个非零投影仍是权 \(\lambda\) 的最高权向量，所以相应简单项均为 \(L(\lambda)\)。若某项投影为零，则整个生成子模在该项的投影为零，与 \(\bar v_\lambda\) 生成全模矛盾。因此所有直和项均为 \(L(\lambda)\)。但 \(V(\lambda)_\lambda\) 一维，而每份 \(L(\lambda)\) 的最高权空间也一维，所以只能有一份，得到同构。
\end{proof}

寻找这个模的一组基时，可以先用
\(\mu(\rho^\vee)\le\lambda(\rho^\vee)\) 枚举可能的支配整权，再取它们的有限 Weyl 轨道，并保留落在 \(\lambda-Q_+\) 中的权；对每个这样的权，PBW 只留下有限多个候选单项式。它们之间仍可能有线性关系，下一章的权重数公式将计算最终维数。

\ProseParagraphBreak
生成元上的局部幂零性要推广到 \(u\bar v_\lambda\)，还须利用 \(\operatorname{ad}e_i,\operatorname{ad}f_i\) 在包络代数中的局部幂零性，以及所引命题的二项式恒等式。有限维性的证明进一步利用整个 Weyl 群的对称性；单个简单根方向的有限权链还不足以限制所有方向的组合。
